\section{Conjuntos conexos y arco-conexos}

\begin{theorem}
  Todo subconjunto no vacío y arco-conexo de $\R^n$ es conexo.
\end{theorem}

\begin{proof}
  Sea $A \subset \R^n$ arco-conexo y no vacío, y sea $a_0 \in A$. Para cada $a \in A$ existe una función continua $\sigma_a\colon [0, 1] \to \R^n$ con valores en $A$ tal que $\sigma_a(0) = a_0$ y $\sigma_a(1) = a$. Como $[0, 1]$ es conexo y $\sigma_a$ es continua, $\sigma_a([0, 1])$ es conexo; además contiene a $a_0$ y a $a$ y está contenido en $A$. Luego $A = \bigcup_{a \in A} \sigma_a([0, 1])$ es unión de conexos que contienen todos a $a_0$, en particular son no disjuntos de a pares, y por lo tanto $A$ es conexo.
\end{proof}

\begin{corollary}
  $\R^n$ es conexo. En particular, los únicos subconjuntos de $\R^n$ que son abiertos y cerrados a la vez son $\varnothing$ y $\R^n$.
\end{corollary}

\begin{proof}
  $\R^n$ es convexo, luego arco-conexo, luego conexo. La segunda afirmación es la caracterización de los conexos por subconjuntos abiertos y cerrados.
\end{proof}

\begin{theorem}
  Un subconjunto abierto y no vacío de $\R^n$ es arco-conexo si y sólo si es conexo.
\end{theorem}

\begin{proof}
  ($\Rightarrow$) Es el teorema anterior.

  ($\Leftarrow$) Sea $A \subset \R^n$ abierto, no vacío y conexo, y sea $x_0 \in A$. Sea $U$ el conjunto de los $x \in A$ para los que existe una función continua $\sigma\colon [0, 1] \to \R^n$ con valores en $A$, $\sigma(0) = x_0$ y $\sigma(1) = x$. Veamos que $U = A$.

  Primero observamos que, si $p \in U$, $q \in A$ y el segmento $\{ (1-s)p + sq : s \in [0, 1] \}$ está contenido en $A$, entonces $q \in U$. En efecto, si $\sigma$ es una curva en $A$ de $x_0$ a $p$, la función
  \[
    \tau(\alpha) = \begin{cases}
      \sigma(2\alpha)                       & \text{si } \alpha \in \left[ 0, \tfrac{1}{2} \right], \\
      (2 - 2\alpha)\, p + (2\alpha - 1)\, q & \text{si } \alpha \in \left[ \tfrac{1}{2}, 1 \right]
    \end{cases}
  \]
  está bien definida, pues ambas expresiones valen $p$ en $\alpha = \frac{1}{2}$, es continua, tiene valores en $A$ y cumple $\tau(0) = x_0$ y $\tau(1) = q$.

  $U$ es abierto: sea $x \in U$. Como $A$ es abierto, existe $r > 0$ con $B_r(x) \subset A$. Si $y \in B_r(x)$, el segmento de $x$ a $y$ está contenido en la bola $B_r(x)$, que es convexa, y por lo tanto en $A$. Luego $y \in U$ y $B_r(x) \subset U$.

  $A \setminus U$ es abierto: sea $u \in A \setminus U$ y $s > 0$ con $B_s(u) \subset A$. Si existiera $y \in B_s(u) \cap U$, el segmento de $y$ a $u$ estaría en $B_s(u) \subset A$ y entonces $u \in U$, absurdo. Luego $B_s(u) \subset A \setminus U$.

  Entonces $A = U \cup (A \setminus U)$ con $U$ y $A \setminus U$ abiertos y disjuntos. Como $A$ es conexo y $x_0 \in U$, debe ser $A \setminus U = \varnothing$, es decir, $U = A$. Esto prueba que $A$ es arco-conexo, pues dados $x, y \in A$ se unen mediante una curva de $x_0$ a $x$ y otra de $x_0$ a $y$; recorriendo la primera al revés y luego la segunda se obtiene una curva de $x$ a $y$ en $A$.
\end{proof}

\section{Continuidad uniforme}

\begin{definition}[Continuidad uniforme]
  Una función $f\colon A \to \R^k$, con $A \subset \R^n$, es \emph{uniformemente continua} si
  \[
    (\forall \e > 0)(\exists \delta > 0) \colon (\forall x, y \in A) \quad \norm{x - y} < \delta \Rightarrow \norm{f(x) - f(y)} < \e.
  \]
\end{definition}

\begin{eg}
  \begin{enumerate}[label=(\arabic*)]
    \item $f(x) = x$ es uniformemente continua: dado $\e > 0$ basta tomar $\delta = \e$.
    \item $f\colon \R \to \R$, $f(x) = \sin x$ es uniformemente continua. Como $f$ es continua en $0$, dado $\e > 0$ existe $\delta' > 0$ tal que $|z| < \delta' \Rightarrow |\sin z| < \frac{\e}{2}$. Sea $\delta = 2\delta'$. Si $|x - y| < \delta$, entonces
          \[
            |\sin x - \sin y| = \left| 2 \sin \frac{x - y}{2} \cos \frac{x + y}{2} \right| \leq 2 \left| \sin \frac{x - y}{2} \right| < 2 \cdot \frac{\e}{2} = \e.
          \]
    \item $f\colon (0, +\infty) \to \R$, $f(x) = \cos \frac{\pi}{x}$ no es uniformemente continua. Sean $x_m = \frac{1}{m}$ e $y_m = \frac{1}{m + 1}$. Entonces $|x_m - y_m| = \frac{1}{m(m+1)} \to 0$, pero $|f(x_m) - f(y_m)| = \bigl| \cos(m\pi) - \cos((m+1)\pi) \bigr| = 2$ para todo $m$. Para $\e = 1$, ningún $\delta > 0$ sirve: existe $m$ con $\frac{1}{m(m+1)} < \delta$ y para ese $m$ se tiene $|x_m - y_m| < \delta$ pero $|f(x_m) - f(y_m)| = 2 \geq \e$.
  \end{enumerate}
\end{eg}

\begin{theorem}
  Si $K \subset \R^n$ es compacto y $f\colon K \to \R^k$ es continua, entonces $f$ es uniformemente continua.
\end{theorem}

\begin{proof}
  Sea $\e > 0$. Como $f$ es continua, para cada $x \in K$ existe $\delta_x > 0$ tal que
  \[
    (\forall z \in K) \quad \norm{z - x} < \delta_x \Rightarrow \norm{f(z) - f(x)} < \frac{\e}{2}.
  \]
  La familia $\mathscr{U} = \{ B_{\delta_x}(x) : x \in K \}$ es un cubrimiento abierto de $K$. Como $K$ es compacto, existe un número de Lebesgue $\delta > 0$ para $\mathscr{U}$: si $x, y \in K$ y $\norm{x - y} < \delta$, existe $z \in K$ con $x, y \in B_{\delta_z}(z)$. Entonces $\norm{f(x) - f(z)} < \frac{\e}{2}$ y $\norm{f(y) - f(z)} < \frac{\e}{2}$, y
  \[
    \norm{f(x) - f(y)} \leq \norm{f(x) - f(z)} + \norm{f(z) - f(y)} < \e.
  \]
  Luego $f$ es uniformemente continua.
\end{proof}

\begin{theorem}[Extensión de funciones uniformemente continuas]
  Sea $A \subset \R^n$ y $f\colon A \to \R^k$ uniformemente continua. Entonces existe una única función continua $\bar{f}\colon \cl{A} \to \R^k$ tal que $\bar{f}(x) = f(x)$ para todo $x \in A$, y además $\bar{f}$ es uniformemente continua.
\end{theorem}

\begin{proof}
  \emph{Paso 1.} Si ${(x_m)} \subset A$ es convergente, entonces ${(f(x_m))}$ es convergente. En efecto, ${(x_m)}$ es de Cauchy. Dado $\e > 0$, sea $\delta > 0$ tal que $\norm{x - y} < \delta \Rightarrow \norm{f(x) - f(y)} < \e$. Existe $m_0$ tal que $\norm{x_m - x_l} < \delta$ para $m, l > m_0$, y entonces $\norm{f(x_m) - f(x_l)} < \e$. Luego ${(f(x_m))}$ es de Cauchy en $\R^k$ y converge.

  \emph{Paso 2.} Si ${(x_m)}, {(y_m)} \subset A$ convergen al mismo punto $x_0$, entonces $\lim f(x_m) = \lim f(y_m)$. En efecto, $\norm{x_m - y_m} \to 0$. Dado $\e > 0$, sea $\delta$ como en la continuidad uniforme; para $m$ grande $\norm{x_m - y_m} < \delta$ y entonces $\norm{f(x_m) - f(y_m)} < \e$. Luego $\norm{f(x_m) - f(y_m)} \to 0$, y los límites, que existen por el paso 1, coinciden.

  \emph{Paso 3: definición de $\bar{f}$.} Dado $x_0 \in \cl{A}$ existe una sucesión ${(x_m)} \subset A$ con $x_m \to x_0$. Definimos $\bar{f}(x_0) = \lim_{m \to \infty} f(x_m)$. Por los pasos 1 y 2, este límite existe y no depende de la sucesión elegida. Si $x \in A$, la sucesión constante $x_m = x$ converge a $x$, de modo que $\bar{f}(x) = f(x)$.

  \emph{Paso 4: $\bar{f}$ es uniformemente continua.} Sea $\e > 0$ y $\delta > 0$ tal que $x, y \in A$, $\norm{x - y} < \delta \Rightarrow \norm{f(x) - f(y)} < \frac{\e}{2}$. Sean $u, v \in \cl{A}$ con $\norm{u - v} < \delta$, y sean ${(u_m)}, {(v_m)} \subset A$ con $u_m \to u$, $v_m \to v$. Entonces $\norm{u_m - v_m} \to \norm{u - v} < \delta$, de modo que $\norm{u_m - v_m} < \delta$ para $m$ grande, y entonces $\norm{f(u_m) - f(v_m)} < \frac{\e}{2}$. Tomando límite cuando $m \to \infty$ se obtiene $\norm{\bar{f}(u) - \bar{f}(v)} \leq \frac{\e}{2} < \e$. En particular, $\bar{f}$ es continua.

  \emph{Paso 5: unicidad.} Sea $g\colon \cl{A} \to \R^k$ continua con $g|_A = f$. Dado $x \in \cl{A}$, sea ${(x_m)} \subset A$ con $x_m \to x$. Por la continuidad de $g$, $g(x_m) \to g(x)$, y por la continuidad de $\bar{f}$, $\bar{f}(x_m) \to \bar{f}(x)$. Como $g(x_m) = f(x_m) = \bar{f}(x_m)$ para todo $m$, por la unicidad del límite $g(x) = \bar{f}(x)$.
\end{proof}

% TODO: Revisar esta sección
\section{Diferenciación}

\begin{definition}
  Sea \(F : A \subset \R^n \to \R^n\) y sea \(p \in A^{\circ} \) si \(T_p : \R^n \to \R^n\) es una transformación lineal que verifica:
  \begin{align*}
    \lim_{x \to p} \dfrac{\| F(x) - F(p) - T_p(x-p) \|}{\|x-p\|} = 0 \Rightarrow
  \end{align*} Decimos que \(F\) es diferenciable en \(p\). A la transformación lineal la llamamos diferencial de \(F\) en \(p\) y lo notamos \(DF_p\).
\end{definition}

Si \(F = (f_1, \ldots, f_n)\) con \(f_i: A \to \R \) \(i \in \{1, \ldots, n\} \), \(F\) es diferenciable en \(p \iff \) cada \(f_i\) es diferenciable en \(p\).
Además la diferencial de \(F\) es la matriz que se obtiene al poner como filas los gradientes de \(f_i\).

\begin{align*}
  DF_p = \begin{pmatrix}
           \nabla f_1(p) \\
           \vdots        \\
           \nabla f_n(p)
         \end{pmatrix} = \begin{pmatrix}
                           \dfrac{\partial f_1}{\partial x_1}(p) & \cdots & \dfrac{\partial f_1}{\partial x_n}(p) \\
                           \vdots                                & \cdots & \vdots                                \\
                           \dfrac{\partial f_n}{\partial x_1}(p) & \cdots & \dfrac{\partial f_n}{\partial x_n}(p)
                         \end{pmatrix}
\end{align*}

También podemos observar que si \(F\) es diferenciable en \(P\) y \(v \in \R^n \Rightarrow \alpha(t) = p + t \cdot v\).

\begin{align*}
  0 & = \lim_{t \to 0} \dfrac{\| F(p + t \cdot v) - F(p) - DF_p(t \cdot v) \|}{|t|}                 \\
    & = \lim_{t \to 0} \| \dfrac{F(p + tv) - F(p)}{t} - DF_p(v) \|                                  \\
    & \iff DF_p(v) = \lim_{t \to 0} \dfrac{F(p+t \cdot v) - F(p)}{t} \text{  } (\forall v \in \R^n)
\end{align*}

\begin{note}
  Si \(G, H: A \subset \R^n \to \R^n\) son diferenciables podemos considerar su producto escalar como \(F(x) = \langle G(x), H(x) \rangle \) que es diferenciable.

  Además:
  \begin{align*}
     & \lim_{t \to 0} \dfrac{F(p + tx) - F(p)}{t} =                                                                                                                 \\
     & \lim_{t \to 0} \dfrac{1}{t} ( \langle G(p + tx), H(p+tx) \rangle - \langle G(p), H(p) \rangle )                                                              \\
     & \lim_{t \to 0} \dfrac{1}{t} ( \langle G(p+tx), H(p+tx) \rangle - \langle G(p), H(p+tx) \rangle + \langle G(p), H(p+tx) \rangle - \langle G(p), H(p) \rangle) \\
     & \lim_{t \to 0} \dfrac{1}{t} (\langle G(p+tx) - G(p), H(p + tx) \rangle + \langle G(p), H(p+tx) - H(p) \rangle)                                               \\
     & \langle DG_p(x), H(p) \rangle + \langle G(p), DH_p \rangle \Rightarrow                                                                                       \\
     & D\langle G, H \rangle_p = \langle DG_p(x), H(p) \rangle + \langle G(p), DH_p(x) \rangle                                                                      \\
  \end{align*}
\end{note}

\begin{lemma}
  Si \(T: \R^n \to \R^n\) es transformación lineal, \(\exists c > 0: \|T(x)\| < c \|x\| \) \((\forall x \in \R^n)\). Al mínimo de estas constantes lo llamamos \( \|T\|_{\infty} \). Vale que \( \|T\|_{\infty} = \max_{\|x\| \leq 1} \|T(x)\| \) en particular
  \( \|T\| \leq \|T\|_{\infty} \|x\| \).
  \begin{proof}
    \(T\) es lineal, luego \(T\) es continua. Tomar norma también es continua así que \(x \mapsto \|Tx\| \) es continua. Como \(\overline{B_1(0)} \) es compacto la función alcanza máximo. Digamos \( \|T\|_{\infty} = \max_{x \in B_r(0)} \|Tx\| \Rightarrow x \neq 0\) \\
    \( \|T(\dfrac{x}{\|x\|})\| \leq \|T\|_{\infty} \Rightarrow \|Tx\| \leq \|x\| \|T\|_{\infty} \) por como la elegimos es la mejor cota.
  \end{proof}
\end{lemma}

\begin{lemma}
  La función \( \| \cdot \|: \R^{n \times m} \to \R \) es una norma.
  \begin{enumerate}
    \item \( \|T\|_{\infty} > 0\), \( \|T\|_{\infty} = 0 \iff T = 0\).
    \item \( \|\lambda \| \|T\|_{\infty} = \|\lambda T\|_{\infty} \).
    \item \( \| T + S\| \leq \|T\|_{\infty} + \|S\|_{\infty} \).
  \end{enumerate}
  \begin{proof}
    \begin{enumerate}
      \item \( \|T\|_{\infty} > 0\) es trivial. Si \(T = 0 \Rightarrow Tx = 0 \Rightarrow \|Tx\|_{\infty} = 0\) y si \( \|T\|_{\infty} \neq 0 \Rightarrow Tx=0 (\forall x \in \overline{B_1(0)}) \) pues por linealidad si
            \begin{align*}
              x \neq 0 \Rightarrow Tx = T(\dfrac{x}{\|x\|} \cdot \|x\|) = \|x\| T(\dfrac{x}{x}) = 0 \therefore T = 0
            \end{align*}
      \item Trivial
      \item \(S + T \in \R^{n \times m} \) dado \(x \in \overline{B_1(0)} \) \begin{align*}
              \|(S+T)(x)\| & = \| Tx + Sx \| \leq \|Tx\| + \|Sx\| \leq  \|x\| \cdot (\|Tx\|_{\infty} + \|Sx\|_{\infty})                   \\
                           & \leq \|Tx\|_{\infty} + \|Sx\|_{\infty} \therefore \|(T+S)(x)\|_{\infty} \leq \|T\|_{\infty} + \|S\|_{\infty}
            \end{align*}
    \end{enumerate}
  \end{proof}
\end{lemma}

\begin{lemma}
  Si \(F: A \subset \R^n \to \R^n\) es diferenciable con \(p \in A^{\circ} \), \(\exists R > 0\) y \(c_r > 0 : x \in B_r(p) \Rightarrow \| F(x) - F(p) \| \leq c_r \| x - p\| \)
  Esto implica la continuidad uniforme en \(B_r(p)\).
  \begin{proof}
    Como \(F\) es diferenciable en \(p\), \(\exists B_r(p)\) donde el cociente de diferenciabilidad es menor que \(1\).\begin{align*}
      \dfrac{\|F(x) - F(p)\|}{\|x - p \|} & \leq \dfrac{\|F(x) - F(p) - DF_p(x-p) \|}{\|x-p\|} + \dfrac{\| DF_p(x-p) \|}{\|x-p\|} \\
                                          & < 1 + \dfrac{\|DF_p(x-p)\|_{\infty} \|x-p\|}{\|x-p\|} = 1 + \|DF_p(x-p)\|_{\infty}
    \end{align*}
    Si \(x \in B_r(p) \Rightarrow \). Luego \(c_r = 1 + \|DF_p(x-p)\|_{\infty} \Rightarrow \|F(X) - F(p)\| \leq c_r \|x-p\| \) \((\forall x \in B_r(p))\).
  \end{proof}
\end{lemma}
